\documentclass[11pt]{article}
\usepackage[a4paper,margin=1in]{geometry}
\usepackage{fontspec}
\usepackage[boldfont=false]{xeCJK}
\setCJKmainfont{FandolSong-Regular.otf}
\usepackage{amsmath}
\usepackage{amssymb}
\usepackage{array}
\usepackage{booktabs}
\usepackage{graphicx}
\usepackage{adjustbox}
\usepackage{enumitem}
\setlist[itemize]{topsep=2pt,partopsep=0pt,parsep=0pt,itemsep=2pt,leftmargin=1.4em}
\setlist[enumerate]{topsep=2pt,partopsep=0pt,parsep=0pt,itemsep=2pt,leftmargin=1.6em}
\usepackage{parskip}
\usepackage{fvextra}
\fvset{breaklines=true,breakanywhere=true,fontsize=\small}
\setlength{\parindent}{0pt}

\begin{document}

\begin{center}
  {\LARGE\bfseries Transformations and vectors}\\[0.35em]
  {\large Igcse Mathematics}
\end{center}
\vspace{0.6em}

This handout covers Topic 7, Transformations and vectors. Parts marked \textbf{(Extended)} are only tested on the Extended papers; everything else is for both levels. Vectors as a whole topic are Extended.

\section*{Transformations}

A \textbf{transformation} 变换 changes the position or size of a shape. The original is the \emph{object} and the result is the \emph{image}. There are four types. When asked to \textbf{describe} one, you must name the type and give all the details below.

\subsection*{Reflection}

A \textbf{reflection} 反射 flips the shape over a \textbf{mirror line} 对称轴. Each image point is the same distance from the line as the object point, on the other side.

\begin{itemize}
\item To describe it, give the \textbf{equation of the mirror line} (for Core, a horizontal or vertical line; Extended allows any line such as $y = x$).

\end{itemize}

\textbf{Worked example.} Reflect the point $(3, 2)$ in the $y$-axis. Only the sign of $x$ changes: the image is $(-3, 2)$. (In the $x$-axis it would be $(3, -2)$.)

\subsection*{Rotation}

A \textbf{rotation} 旋转 turns the shape about a fixed point, the \textbf{centre} 中心 of rotation, through multiples of $90^{\circ}$.

\begin{itemize}
\item To describe it, give the centre, the angle, and the direction (clockwise or anticlockwise).

\end{itemize}

\textbf{Worked example.} Rotate $(3, 1)$ by $90^{\circ}$ anticlockwise about the origin. The rule is $(x, y) \to (-y, x)$, so the image is $(-1, 3)$.

\subsection*{Enlargement}

An \textbf{enlargement} 放大 changes the size by a \textbf{scale factor} 比例因子 $k$, measured from a fixed centre. Each distance from the centre is multiplied by $k$.

\begin{itemize}
\item To describe it, give the centre and the scale factor. For Extended, $k$ may be \textbf{fractional} (the shape gets smaller) or \textbf{negative} (the image appears on the other side of the centre).

\end{itemize}

\textbf{Worked example.} Enlarge $(1, 2)$ from the origin by scale factor $2$. Multiply both coordinates: the image is $(2, 4)$.

\subsection*{Translation}

A \textbf{translation} 平移 slides the shape with no turning, by a \textbf{vector} 向量 written as a \textbf{column vector} 列向量 $\begin{pmatrix} x \\ y \end{pmatrix}$ ($x$ across, $y$ up).

\textbf{Worked example.} Translate $(5, 3)$ by $\begin{pmatrix} -2 \\ 4 \end{pmatrix}$: move $2$ left and $4$ up to get $(3, 7)$.

(\textbf{Extended:} a question may ask you to combine two transformations and describe the single transformation that has the same effect.)

\section*{Vectors in two dimensions (Extended)}

A vector has both size and direction. It can be written as a column vector, as $\overrightarrow{AB}$ (from $A$ to $B$), or in bold as $\mathbf{a}$.

\begin{itemize}
\item \textbf{Add or subtract} by working on the top and bottom numbers separately.

\item \textbf{Multiply by a scalar} 标量 (an ordinary number) by multiplying both numbers.

\end{itemize}

\textbf{Worked example.} If $\mathbf{a} = \begin{pmatrix} 3 \\ 1 \end{pmatrix}$ and $\mathbf{b} = \begin{pmatrix} 2 \\ -4 \end{pmatrix}$, then

\[
\mathbf{a} + \mathbf{b} = \begin{pmatrix} 5 \\ -3 \end{pmatrix}, \qquad 3\mathbf{a} = \begin{pmatrix} 9 \\ 3 \end{pmatrix}.
\]

\section*{Magnitude of a vector (Extended)}

The \textbf{magnitude} 模 (length) of a vector is found with Pythagoras. For $\begin{pmatrix} x \\ y \end{pmatrix}$,

\[
\left| \begin{pmatrix} x \\ y \end{pmatrix} \right| = \sqrt{x^{2} + y^{2}}.
\]

\textbf{Worked example.} The magnitude of $\begin{pmatrix} 3 \\ 4 \end{pmatrix}$ is $\sqrt{3^{2} + 4^{2}} = \sqrt{25} = 5$.

\section*{Vector geometry (Extended)}

A vector can be drawn as a \textbf{directed line segment} 有向线段 (an arrow). The \textbf{position vector} 位置向量 of a point is the vector from the origin $O$ to that point.

A key idea: the vector from $A$ to $B$ is

\[
\overrightarrow{AB} = \mathbf{b} - \mathbf{a},
\]

where $\mathbf{a}$ and $\mathbf{b}$ are the position vectors of $A$ and $B$.

Two vectors are \textbf{parallel} 平行 if one is a scalar multiple of the other (for example $\overrightarrow{AB} = 2\,\overrightarrow{CD}$). Three points are \textbf{collinear} 共线 (in a straight line) if the vectors between them are parallel and share a point. You can express any vector in terms of two \textbf{coplanar} 共面 vectors.

\textbf{Worked example.} $O$ is the origin, with $\overrightarrow{OA} = \mathbf{a}$ and $\overrightarrow{OB} = \mathbf{b}$. $M$ is the midpoint of $AB$. Find $\overrightarrow{OM}$ in terms of $\mathbf{a}$ and $\mathbf{b}$.

\[
\overrightarrow{OM} = \mathbf{a} + \tfrac{1}{2}\overrightarrow{AB} = \mathbf{a} + \tfrac{1}{2}(\mathbf{b} - \mathbf{a}) = \tfrac{1}{2}(\mathbf{a} + \mathbf{b}).
\]


\end{document}
